\chapter{Dynamics Processing}
\label{chap:dynamics}

Dynamics processing changes not frequency but \textbf{how level behaves over time}. It makes loud peaks step back, brings quiet detail closer, or keeps unwanted sound out. Compression, limiting, gating, and expansion are four different goals; treating "loud" as the only goal usually flattens a mix.

\section{Four devices, four goals}

\begin{table}[htbp]
\centering
\small
\begin{tabular}{@{}p{0.16\linewidth}p{0.36\linewidth}p{0.38\linewidth}@{}}
\hline
Device & What it does & Typical use \\
\hline
Compressor & Reduces the part above threshold by a ratio; below threshold stays untouched & Control vocal/bass dynamics; make a track sit steadily in the mix \\
Limiter & Pins peaks below a ceiling with a very high ratio and very fast attack & Prevent overs before export, or protect downstream devices \\
Gate & Attenuates or mutes the part below threshold & Remove noise floor, bleed, unwanted tails \\
Expander & Pushes the part below threshold further down; upward expansion lifts it & Restore contrast, make quieter sections quieter; upward expansion is rarer \\
\hline
\end{tabular}
\caption{The different goals of four dynamics devices}
\label{tab:dynamics-kinds}
\end{table}

The line between compression and limiting is not the name but \textbf{whether the attack lets transients through}. A compressor can use a slow attack so the first transient stays intact and only the sustained part is compressed; a limiter's attack must be fast enough that the peak does not cross the ceiling. A gate is not a "reverse compressor": it acts below threshold to remove content. An expander also acts below threshold, but its goal is to widen the level difference, not to mute.

\section{Compressor parameters}

Yeban's built-in device is a feed-forward soft-knee compressor. The chain is: stereo-linked RMS detection → soft-knee static curve → independent gain attack/release → multiply back into the input, then makeup. It has no lookahead delay and does not look at the output, so there is no feedback loop.

\begin{center}
\begin{tikzpicture}[font=\scriptsize, node distance=5mm and 8mm,
  box/.style={draw, rounded corners=2pt, minimum width=24mm, minimum height=8mm, align=center, fill=black!5},
  arrow/.style={-{Latex[length=1.8mm]}, semithick}]
  \node[box] (in) {Input L/R};
  \node[box, right=of in] (det) {Linked RMS\\detector};
  \node[box, right=of det] (curve) {Soft-knee\\static curve};
  \node[box, right=of curve] (ball) {Attack/release\\gain ballistics};
  \node[box, right=of ball] (out) {Multiply back\\+ makeup};
  \draw[arrow] (in) -- (det);
  \draw[arrow] (det) -- (curve);
  \draw[arrow] (curve) -- (ball);
  \draw[arrow] (ball) -- (out);
\end{tikzpicture}
\end{center}

\begin{table}[htbp]
\centering
\small
\begin{tabular}{@{}llll@{}}
\hline
Parameter & Meaning & Range & Default \\
\hline
\texttt{threshold\_db} & Where compression starts & $-120$–$24$ dBFS & $-12$ dBFS \\
\texttt{ratio} & Compression ratio & 1–1000 & 4:1 \\
\texttt{knee\_db} & Soft-knee width & 0–24 dB & 6 dB \\
\texttt{detector\_s} & RMS detector window time constant & 0.0001–10 s & 0.005 s \\
\texttt{attack\_s} & Time constant when gain moves down & 0.0001–10 s & 0.01 s \\
\texttt{release\_s} & Time constant for gain returning to 1.0 & 0.0001–10 s & 0.1 s \\
\texttt{makeup\_db} & Post-compression makeup gain & $-24$–$24$ dB & 0 dB \\
\hline
\end{tabular}
\caption{Yeban's dynamics parameter contract}
\label{tab:comp-params}
\end{table}

All three time values are \textbf{time constants $\tau$}, not "10%–90% rise time": 63.2\% of the way to steady state is $\tau$, and 1 dB from steady state is about $1.585\tau$. The detector window answers "how fast do we see the current RMS"; attack/release answer "how fast does gain follow". Merging them into one attack knob makes the readings impossible to reconcile.

The soft-knee formula interpolates quadratically within $W/2$ above and below threshold: transparent below the knee, smooth inside, ratio compression above. $W=0$ degenerates to a hard knee with a corner at threshold. The point of a soft knee is to avoid abruptly crushing material that just crosses threshold, not to make it louder.

\section{Where the compressor actually lives in Yeban}

Click "Insert Compressor" in the device rack to add the built-in soft-knee compressor to the current track. The rack's compressor readout shows threshold, ratio, and GR; the button beside it cycles the ratio: \textbf{2:1 → 4:1 → 8:1 → 20:1}. If the current ratio is not in that list, the next click starts at 2:1.

\begin{table}[htbp]
\centering
\small
\begin{tabular}{@{}p{0.28\linewidth}p{0.34\linewidth}p{0.30\linewidth}@{}}
\hline
UI element & Actual behavior & Boundary \\
\hline
"Insert Compressor" & Inserts the dynamics group with contract defaults & Multiple units can be inserted on one track \\
Compressor readout & Shows threshold/ratio of the first device with \texttt{ratio} & Falls back to placeholder values when none is found \\
Ratio cycle button & Changes only that device's \texttt{ratio} & Only the four values 2/4/8/20 \\
GR readout & Always shows 0.0 dB & Not a live gain-reduction meter \\
Compressor curve & Draws a hard-knee transfer curve from threshold + ratio & Ignores knee/detector/attack/release; not a soft-knee measurement \\
Threshold/knee/detector/attack/release/makeup knobs & The contract has the parameters & The current rack has no separate controls; they cannot be dragged \\
\hline
\end{tabular}
\caption{What the compressor UI elements actually do}
\label{tab:comp-ui}
\end{table}

The two easiest misreads are "GR 0.0 dB" and the "measured curve". The GR string is fixed by the projection and does not mean there is no reduction; the curve uses the hard-knee formula $y=T+(x-T)/R$, while the device runs the three-part soft-knee formula. Judge whether compression happens by ear, meters, and the overall dynamics; do not treat the curve as a live meter.

\screenshot{device-rack-compressor}{The current track has a Compressor inserted; the device rack shows the compressor readout and Ratio cycle button, with the compressor curve visible.}

\section{Status of limiter, gate, and expander}

Yeban's effect contract has a separate limiter group: threshold $-30$–$0$ dB (default $-1$), ceiling $-6$–$0$ dB (default $-0.3$), release 10–1000 ms (default 80), lookahead 1–192 frames. But the one-click built-in effects are only EQ, Reverb, and Compressor; there is no "Insert Limiter" entry. The UI line "Limiter: threshold -0.1 dBTP · status normal (no overs)" is fixed text, not a state read from a real limiter device.

Gate and expander have no corresponding built-in effect group and no rack entry. They can be explained as concepts, but they cannot be operated as existing features on this rack. If these processes are needed, use a path outside the current built-in contracts or wait for a later version; the manual does not write "should exist" as "exists now".

\section{Setup order and common mistakes}

\begin{enumerate}
  \item \textbf{Fix level before compression}: if a track jumps because of performance, flatten the sections with fader/automation first; compression handles the remaining peaks and dips.
  \item \textbf{Set threshold against the loudest section}: it should sit below the peaks that need control without compressing every phrase; a common start is 3–6 dB reduction on the loudest part.
  \item \textbf{Choose ratio by goal}: 2:1–4:1 suits vocal/bass dynamics; 8:1 and above feels closer to limiting and can sound squashed.
  \item \textbf{Do not always use the fastest attack}: too fast shaves drum/plucked transients; too slow lets peaks through. 10 ms is a common start.
  \item \textbf{Do not always use the shortest release}: too short pumps; too long makes the next phrase slow to recover. 100 ms is a common start, adjusted to the song.
  \item \textbf{Makeup is not a loudness button}: it raises everything and can clip the bus; check post-compression peaks before adding gain.
\end{enumerate}

Common mistakes: using a compressor as a limiter (slow attack leaks peaks), using a limiter as a loudness tool (more limiting means flatter), using a gate to "fix" performance (it cuts tails and breaths too), trusting the fixed GR 0.0 as "no compression", and reading the hard-knee schematic as a soft-knee measurement. Yeban currently has no separate knobs for these parameters; when fine control is needed, solve it in project data or wait for later controls, and do not pretend curve dragging changes knee or attack.

\section{Exercise}

\begin{enumerate}
  \item Put a compressor on a track with wide dynamics. With the threshold fixed, compare 2:1 and 8:1; note transients and sustained parts.
  \item Treat the GR reading as a placeholder, not proof of ``no compression''; judge by ear, meters, and overall movement.
  \item Do a ``fix level first, then compress'' A/B: flatten sections with fader/automation, then compress; compare pumping and detail.
\end{enumerate}

\section{Help-topic mapping}

\begin{table}[htbp]
\centering
\small
\begin{tabular}{@{}p{0.40\linewidth}p{0.48\linewidth}@{}}
\hline
Section & In-app help anchor \\
\hline
Four devices, four goals & \texttt{dyn-four-kinds} \\
Compressor parameters & \texttt{dyn-compressor-params} \\
Where the compressor actually lives in Yeban & \texttt{dyn-compressor-ui} \\
Status of limiter, gate, and expander & \texttt{dyn-limiter-gate-status} \\
Setup order and common mistakes & \texttt{dyn-workflow-mistakes} \\
\hline
\end{tabular}
\caption{Help-topic mapping for Chapter 4}
\end{table}
